% !TeX root = ../mathematical-methods1.tex
\chapter{Differential Calculus}\label{chap:differential-calculus}

This chapter concerns functions of one real variable. It follows the lecturer's
six topics, including the introduction to integration and series approximations.
The examples retain their lecture order; selected lecture diagrams are redrawn
schematically. Conditions that were implicit in the manuscript are stated explicitly.
Throughout this chapter, $\log$ means the natural logarithm and angles are in radians.

\section{The derivative}

\begin{definition}[Differentiability]
Let $(a,b)\subseteq\mathbb R$ and $f:(a,b)\to\mathbb R$.
The function $f$ is differentiable at $x\in(a,b)$ if the finite limit
\begin{equation}
 f'(x)=\lim_{h\to0}\frac{f(x+h)-f(x)}{h}
 \tag{1}\label{eq:mm-derivative}
\end{equation}
exists. Other notation is $\mathrm df/\mathrm dx$.
The function is differentiable on $(a,b)$ if this derivative exists at every
point of the interval.
\end{definition}

Geometrically, $f'(x)$ is the slope of the tangent to the graph at $x$.
The difference quotient is the slope of the secant through the points with
horizontal coordinates $x$ and $x+h$; the tangent slope is its limit.
In applications a derivative measures an instantaneous rate of change.
For example, if $x(t)$ is position at time $t$, then $x'(t)$ is velocity.
A rigorous treatment of differentiation is developed in Analysis~1.

\begin{center}
\begin{tikzpicture}[x=1.1cm,y=0.7cm,>=Stealth]
 \draw[->] (-0.2,0)--(4.4,0) node[right] {$x$};
 \draw[->] (0,-0.15)--(0,4.4) node[above] {$f$};
 \draw[thick,ink,domain=0.15:3.8,samples=60] plot (\x,{0.22*\x*\x+0.45});
 \draw[muted,domain=0.4:3.9] plot (\x,{0.88*\x-0.43});
 \draw[blue,domain=0.6:3.9] plot (\x,{1.144*\x-0.958});
 \fill (2,1.33) circle (2pt) node[above left] {$P$};
 \fill (3.2,2.7028) circle (2pt) node[above left] {$Q$};
 \draw[dashed] (2,0) node[below] {$x$}--(2,1.33);
 \draw[dashed] (3.2,0) node[below] {$x+h$}--(3.2,2.7028);
 \node[anchor=west,font=\small] at (4.6,3.2) {Secant: $P,Q$};
 \node[anchor=west,font=\small] at (4.6,2.4) {Tangent: limit as $h\to0$};
\end{tikzpicture}
\end{center}

\begin{example}[Lecture Example 1: use the definition]
Find the derivative of $f(x)=x^2$ directly from \eqref{eq:mm-derivative}.
\begin{align*}
 f'(x)
 &=\lim_{h\to0}\frac{(x+h)^2-x^2}{h}\\
 &=\lim_{h\to0}\frac{x^2+2xh+h^2-x^2}{h}\\
 &=\lim_{h\to0}(2x+h)=2x.
\end{align*}
\end{example}

\subsection{Basic derivative formulas}
\begin{center}
\begin{tabular}{lll}
\toprule
$f(x)$ & $f'(x)$ & Domain used here\\
\midrule
$x^\alpha$ & $\alpha x^{\alpha-1}$ & $x>0$, $\alpha\in\mathbb R$\\
$e^x$ & $e^x$ & $x\in\mathbb R$\\
$\log x$ & $1/x$ & $x>0$\\
$\sin x$ & $\cos x$ & $x\in\mathbb R$\\
$\cos x$ & $-\sin x$ & $x\in\mathbb R$\\
\bottomrule
\end{tabular}
\end{center}
For integer powers the power rule also applies at other points where the
function and the displayed derivative are defined.

\subsection{Rules of differentiation}
If $f$ and $g$ are differentiable and $\lambda$ is constant, then
\begin{align*}
 (f+g)'&=f'+g', & (\lambda f)'&=\lambda f',\\
 (fg)'&=f'g+fg'.
\end{align*}
The first two identities express linearity of differentiation.
For a composition, the chain rule is
\[
 \frac{\mathrm d}{\mathrm dx}g(f(x))=g'(f(x))f'(x),
 \qquad (g\circ f)'=(g'\circ f)f'.
\]
These rules follow from the derivative definition using the limit tools of analysis.

\section{Inverse functions and their derivatives}

If $f$ is strictly monotone on an interval, it is injective there. In the
increasing case, $x_1>x_2$ implies $f(x_1)>f(x_2)$; in the decreasing case,
$x_1>x_2$ implies $f(x_1)<f(x_2)$.
Taking the codomain to be its image, $f$ therefore has an inverse:
\[
 y=f(x)\quad\Longleftrightarrow\quad x=f^{-1}(y).
\]
Write $g=f^{-1}$. Then $g(f(x))=x$.
For a differentiable, strictly monotone $f$, at points where $f'(x)\ne0$,
the inverse is differentiable and the chain rule gives
\[
 g'(f(x))f'(x)=1.
\]
Substituting $x=f^{-1}(y)$ yields
\begin{equation}
 g'(y)=\frac{1}{f'(f^{-1}(y))}.
 \tag{2}\label{eq:mm-inverse}
\end{equation}
In function notation,
\begin{equation}
 (f^{-1})'=\frac{1}{f'\circ f^{-1}}.
 \tag{2'}
\end{equation}
\begin{remark}
The nonzero-derivative condition is essential for this finite derivative formula.
The notation $f^{-1}$ here means an inverse function, not $1/f$.
\end{remark}

\begin{example}[Lecture Example 2: the inverse sine]
Find the derivative of $\arcsin x=\sin^{-1}x$ for $-1<x<1$.
Restrict sine to $(-\pi/2,\pi/2)$, on which it is strictly increasing and has
image $(-1,1)$. Formula \eqref{eq:mm-inverse} gives
\[
 (\arcsin x)'=\frac{1}{\cos(\arcsin x)}.
\]
Put $\theta=\arcsin x$. Then $\sin\theta=x$ and
$-\pi/2<\theta<\pi/2$, so $\cos\theta>0$. Consequently,
\[
 \cos(\arcsin x)=\sqrt{1-\sin^2\theta}=\sqrt{1-x^2},
 \qquad \boxed{(\arcsin x)'=\frac{1}{\sqrt{1-x^2}}}.
\]
The positive square root follows from the chosen range of $\arcsin$.
\end{example}

\begin{center}
\begin{tikzpicture}[scale=1.05,>=Stealth]
 \draw[->] (-1.9,0)--(1.9,0) node[right] {$x$};
 \draw[->] (0,-1.3)--(0,1.4) node[above] {$y$};
 \draw[ink,thick,domain=-1.5708:1.5708,samples=60]
   plot (\x,{sin(deg(\x))});
 \node[below] at (-1.57,0) {$-\pi/2$};
 \node[below] at (1.57,0) {$\pi/2$};
 \node[above,font=\small] at (0,1.65) {$y=\sin x$, restricted domain};
 \begin{scope}[xshift=6cm]
  \draw[->] (-1.4,0)--(1.5,0) node[right] {$x$};
  \draw[->] (0,-1.85)--(0,1.9) node[above] {$y$};
  \draw[ink,thick,domain=-1.5708:1.5708,samples=60]
    plot ({sin(deg(\x))},\x);
  \node[below] at (-1,0) {$-1$};
  \node[below] at (1,0) {$1$};
  \node[above,font=\small] at (0,2.15) {$y=\arcsin x$};
 \end{scope}
\end{tikzpicture}
\end{center}

\section{Implicit differentiation}

The inverse-function discussion concerns cases where one variable can be
expressed in terms of the other. More generally, a relation between $x$ and $y$
may define a curve without a convenient explicit formula $y=f(x)$.
To find its slope locally, regard $y$ as a differentiable function $y(x)$
and differentiate the relation with respect to $x$. Apply the chain rule
whenever differentiating an expression involving $y(x)$.

\begin{example}[Lecture Example 3: slope of an implicit curve]
Find the slope at $(x,y)=(2,1)$ of the curve
\[
 y^7+x^4y^3-10x+3=0.
\]
First, $1+16-20+3=0$, so the point lies on the curve.
Differentiating with respect to $x$ gives
\begin{align*}
 0
 &=\frac{\mathrm d}{\mathrm dx}
   \bigl(y(x)^7+x^4y(x)^3-10x+3\bigr)\\
 &=7y^6y'+4x^3y^3+3x^4y^2y'-10.
\end{align*}
The first term uses the chain rule. Differentiating $x^4y(x)^3$ uses both
the product rule and the chain rule. Collect the terms involving $y'$:
\begin{align*}
 y'(7y^6+3x^4y^2)&=10-4x^3y^3,\\
 y'&=\frac{10-4x^3y^3}{7y^6+3x^4y^2}.
\end{align*}
At $(2,1)$ the denominator is $55\ne0$, and
\[
 \left.\frac{\mathrm dy}{\mathrm dx}\right|_{(2,1)}
 =\frac{10-4\cdot2^3\cdot1^3}{7\cdot1^6+3\cdot2^4\cdot1^2}
 =-\frac{22}{55}=-\frac25.
\]
\end{example}

\begin{example}[Lecture Example 4: a self-exponential]
For $x>0$, find the derivative of $f(x)=x^x$.
Let $y=x^x$. Taking natural logarithms gives $\log y=x\log x$.
Differentiating this identity,
\begin{align*}
 \frac{y'}{y}&=\log x+x\frac1x=\log x+1,\\
 y'&=y(\log x+1).
\end{align*}
Therefore
\[
 \boxed{f'(x)=x^x(\log x+1)}.
\]
\end{example}

\section{Applications of derivatives}

\subsection{Optimisation: maxima and minima}
\begin{theorem}[Necessary condition for an interior extremum]
Let $f$ be differentiable on $(a,b)$ and let $c\in(a,b)$.
If $f$ has a local maximum or minimum at $c$, then $f'(c)=0$.
\end{theorem}

\begin{proof}[Proof sketch from the lecture]
Suppose $f$ has a local maximum at $c$. For sufficiently small $h$,
$f(c+h)-f(c)\le0$. Hence
\[
 \frac{f(c+h)-f(c)}{h}\le0\quad(h>0),\qquad
 \frac{f(c+h)-f(c)}{h}\ge0\quad(h<0).
\]
Differentiability means that the two one-sided limits exist and agree.
The first inequality implies $f'(c)\le0$ and the second implies
$f'(c)\ge0$. Thus $f'(c)=0$. The minimum case is analogous.
\end{proof}

To find an absolute maximum or minimum of a differentiable function on a closed
interval, find interior stationary points and compare the values there with
the endpoint values. For continuous functions that fail to be differentiable
at some interior points, those points must also be checked.
A stationary point alone does not guarantee an extremum.

\begin{example}[Lecture Example 5: include the endpoints]
Find the maximum and minimum of $f(x)=x-2\sin x$ on $[0,2\pi]$.
\begin{align*}
 f'(x)&=1-2\cos x,\\
 f'(x)=0&\quad\Longleftrightarrow\quad\cos x=\frac12.
\end{align*}
The interior candidates are $x=\pi/3$ and $x=5\pi/3$.
Their values and the endpoint values are
\begin{align*}
 f(\pi/3)&=\frac\pi3-\sqrt3\approx-0.685,\\
 f(5\pi/3)&=\frac{5\pi}3+\sqrt3\approx6.968,\\
 f(0)&=0, & f(2\pi)&=2\pi\approx6.283.
\end{align*}
Thus the absolute minimum is $\pi/3-\sqrt3$ at $x=\pi/3$,
and the absolute maximum is $5\pi/3+\sqrt3$ at $x=5\pi/3$.
\end{example}

\begin{center}
\begin{tikzpicture}[x=1.05cm,y=0.38cm,>=Stealth]
 \draw[->] (-0.3,0)--(6.8,0) node[right] {$x$};
 \draw[->] (0,-1.8)--(0,8.8) node[above] {$f(x)$};
 \draw[ink,thick,domain=0:6.283185,samples=100]
   plot (\x,{\x-2*sin(deg(\x))});
 \fill (1.047198,-0.68485) circle (2pt);
 \fill (5.235988,6.96804) circle (2pt);
 \draw[dashed] (5.235988,0)--(5.235988,6.96804);
 \node[below] at (1.047198,-0.8) {$\pi/3$};
 \node[below] at (5.235988,0) {$5\pi/3$};
 \node[below] at (6.283185,0) {$2\pi$};
 \node[above] at (5.235988,6.96804) {maximum};
 \node[below] at (1.047198,-1.8) {minimum};
\end{tikzpicture}
\end{center}

\subsection{First-order approximation}
The tangent line to $f$ at $x=c$ has equation
\[
 f_{\mathrm t}(x)-f(c)=f'(c)(x-c).
\]
For $x$ close to $c$, this gives the first-order approximation
\[
 \boxed{f(x)\approx f(c)+f'(c)(x-c)}.
\]
\begin{example}[Lecture Example 6: estimate $\sqrt{17}$]
Take $f(x)=\sqrt{x}$ and $c=16$. Then
\[
 f(16)=4,\qquad f'(x)=\frac{1}{2\sqrt{x}},\qquad f'(16)=\frac18.
\]
Consequently,
\[
 \sqrt{17}=f(17)\approx4+\frac18(17-16)=4.125.
\]
A calculator gives $\sqrt{17}\approx4.123106$, so the absolute error
is about $0.0019$, or approximately $0.002$ as noted in the lecture.
\end{example}

\section{Integration}

Geometrically, a definite integral gives signed area under a graph.
In applications it represents accumulation of a changing quantity: integrating
velocity gives displacement; integrating speed gives distance travelled.
The lecture introduces integration through areas of piecewise constant
approximations; a rigorous treatment belongs to Analysis~2.

\subsection{Partitions and sums of rectangle areas}
Let $f:[a,b]\to\mathbb R$ be continuous and take a partition
\[
 a=x_0<x_1<\cdots<x_{n-1}<x_n=b.
\]
A right-endpoint rectangle on $[x_{i-1},x_i]$ has signed area
$f(x_i)(x_i-x_{i-1})$. Adding the areas gives a Riemann sum.
For equal-width intervals, put
\[
 \Delta x=\frac{b-a}{n},\qquad x_i=a+i\Delta x.
\]
The definite integral is then
\[
 \boxed{\int_a^b f(x)\,\mathrm dx
 =\lim_{n\to\infty}\sum_{i=1}^n f(x_i)\Delta x}.
\]
For arbitrary partitions, the limiting requirement is that the largest
subinterval width tends to zero; merely increasing the number of intervals
is not enough. The continuity assumption ensures the sums converge to the
same value for these choices.

\begin{example}[Lecture Example 7: integrate directly from the definition]
For $f(t)=t$ and $x>0$, find $\int_0^x f(t)\,\mathrm dt$.
Divide $[0,x]$ into $n$ equal intervals with endpoints $t_i=ix/n$.
\begin{align*}
 \int_0^x t\,\mathrm dt
 &=\lim_{n\to\infty}\sum_{i=1}^n
    \left(\frac{ix}{n}\right)\frac{x}{n}\\
 &=\lim_{n\to\infty}\frac{x^2}{n^2}\sum_{i=1}^n i\\
 &=\lim_{n\to\infty}\frac{x^2}{n^2}\frac{n(n+1)}2\\
 &=x^2\lim_{n\to\infty}\frac{n(n+1)}{2n^2}
 =\frac{x^2}{2}.
\end{align*}
The identity $\sum_{i=1}^n i=n(n+1)/2$ is used in the third line.
\end{example}

\begin{center}
\begin{tikzpicture}[scale=0.9,>=Stealth]
 \foreach \i in {1,...,6} {
  \pgfmathsetmacro{\l}{(\i-1)*0.5}
  \pgfmathsetmacro{\r}{\i*0.5}
  \draw[draw=ink,fill=ink!8] (\l,0) rectangle (\r,\r);
 }
 \draw[->] (-0.2,0)--(3.7,0) node[right] {$t$};
 \draw[->] (0,-0.2)--(0,3.7) node[above] {$f(t)$};
 \draw[thick,blue] (0,0)--(3.3,3.3) node[right] {$f(t)=t$};
 \node[below] at (0,0) {$0$};
 \node[below] at (3,0) {$x$};
 \node[below,font=\small] at (1.5,-0.65) {Right-endpoint rectangles};
\end{tikzpicture}
\end{center}

\subsection{The fundamental theorem of calculus}
In the preceding example,
\[
 F(x)=\int_0^x t\,\mathrm dt=\frac{x^2}{2}
 \quad\text{has}\quad F'(x)=x=f(x).
\]
This relation holds more generally.
\begin{theorem}[\textbf{Fundamental theorem of calculus}]
Let $f$ be continuous on $[a,b]$ and set
\[
 F(x)=\int_a^x f(t)\,\mathrm dt.
\]
Then $F'(x)=f(x)$ for $a<x<b$.
Moreover, if $G$ is any antiderivative of $f$, so $G'=f$, then
\[
 \boxed{\int_a^b f(t)\,\mathrm dt=G(b)-G(a)}.
\]
This result is also called the Newton--Leibniz theorem.
\end{theorem}

\textbf{Geometric proof idea.}
For small $h$, additivity of area gives
\[
 F(x+h)=F(x)+\int_x^{x+h}f(t)\,\mathrm dt.
\]
Continuity makes the narrow strip approximately a rectangle of height $f(x)$
and signed width $h$, so
\[
 F(x+h)-F(x)\approx f(x)h,
 \qquad \frac{F(x+h)-F(x)}h\longrightarrow f(x).
\]
This explains why differentiating accumulated area recovers the original
function. The approximation here is a proof idea, rather than a replacement
for the limit argument in Analysis~2.

\subsection{Antiderivatives and indefinite integrals}
An antiderivative (or primitive) of $f$ is a function $G$ satisfying $G'=f$.
The notation
\[
 \int f(x)\,\mathrm dx=G(x)+C
\]
represents the family of antiderivatives, which differ by an additive constant
on an interval. For example,
\[
 \int x\,\mathrm dx=\frac{x^2}{2}+C.
\]
A definite integral is a number; an indefinite integral describes functions.
The fundamental theorem explains the sense in which integration reverses
differentiation.

\subsection{Properties and elementary formulas}
For integrable $f,g$ and a constant $\lambda$,
\begin{align*}
 \int_a^b(f+g)\,\mathrm dx
   &=\int_a^b f\,\mathrm dx+\int_a^b g\,\mathrm dx,\\
 \int_a^b\lambda f\,\mathrm dx
   &=\lambda\int_a^b f\,\mathrm dx,\\
 \int_a^b f\,\mathrm dx
   &=\int_a^c f\,\mathrm dx+\int_c^b f\,\mathrm dx
      \qquad(a<c<b).
\end{align*}
The first two properties express linearity of integration.

\begin{center}
\begin{tabular}{lll}
\toprule
$f(x)$ & $\displaystyle\int f(x)\,\mathrm dx$ & Condition\\[5pt]
\midrule
$x^\alpha$ & $\displaystyle\frac{x^{\alpha+1}}{\alpha+1}+C$
 & $x>0$, $\alpha\ne-1$\\[9pt]
$e^x$ & $e^x+C$ & $x\in\mathbb R$\\[4pt]
$1/x$ & $\log|x|+C$ & an interval avoiding $0$\\[4pt]
$\sin x$ & $-\cos x+C$ & $x\in\mathbb R$\\[4pt]
$\cos x$ & $\sin x+C$ & $x\in\mathbb R$\\
\bottomrule
\end{tabular}
\end{center}
For a definite integral the same constant cancels:
$(G(b)+C)-(G(a)+C)=G(b)-G(a)$.

\section{Series approximation}

\subsection{From a tangent line to a quadratic}
Recall the first-order approximation near $x=c$:
\[
 f(x)\approx f(c)+f'(c)(x-c).
\]
To understand higher orders, first take the centre to be $0$.
The linear polynomial $P_1(x)=f(0)+f'(0)x$ matches both the value and
the first derivative of $f$ at $0$.
For a second-order approximation, try
\[
 P_2(x)=f(0)+f'(0)x+Ax^2.
\]
Matching derivatives at $0$ gives
\begin{align*}
 P_2'(x)&=f'(0)+2Ax,& P_2'(0)&=f'(0),\\
 P_2''(x)&=2A,& P_2''(0)&=2A=f''(0).
\end{align*}
Thus $A=f''(0)/2$, and
\[
 f(x)\approx f(0)+f'(0)x+\frac{f''(0)}2x^2.
\]

\textbf{Exercise from handwritten p.~60.}
Find a second-order approximation to $\sqrt{17}$ and compare it with a
calculator. The margin of the lecture records approximately $4.123046$
for the quadratic approximation, compared with $4.123106$ from a calculator.
% Your own derivation: use f(x)=sqrt(x) centred at 16, not at 0.
% The exact quadratic estimate is 4 + 1/8 - 1/512 = 4.123046875.

\subsection{Maclaurin and Taylor polynomials}
Matching further derivatives extends the construction to degree $N$:
\[
 P_{N,0}(x)=f(0)+f'(0)x+\frac{f''(0)}{2!}x^2
 +\frac{f'''(0)}{3!}x^3+\cdots+\frac{f^{(N)}(0)}{N!}x^N.
\]
The associated infinite series centred at $0$ is the \emph{Maclaurin series}:
\begin{equation}
 \sum_{n=0}^{\infty}\frac{f^{(n)}(0)}{n!}x^n.
 \tag{3}\label{eq:mm-maclaurin}
\end{equation}
For a general centre $a$, the degree-$N$ Taylor polynomial is
\[
 P_{N,a}(x)=\sum_{n=0}^N\frac{f^{(n)}(a)}{n!}(x-a)^n,
\]
and the associated \emph{Taylor series} is
\begin{equation}
 \sum_{n=0}^{\infty}\frac{f^{(n)}(a)}{n!}(x-a)^n.
 \tag{3'}\label{eq:mm-taylor}
\end{equation}
Here $f^{(0)}=f$ and $0!=1$.

\begin{remark}[Approximation versus equality]
The lecture introduces these expressions using an approximation sign.
A finite Taylor polynomial approximates $f$ near its centre under suitable
smoothness assumptions. An infinite Taylor series equals $f(x)$ only where
the remainder $f(x)-P_{N,a}(x)$ tends to zero as $N\to\infty$.
The existence of derivatives of every order alone does not ensure this.
This distinction will be treated more fully in analysis.
\end{remark}

\begin{example}[Lecture Example 8: the exponential function]
Let $f(x)=e^x$. Every derivative is $e^x$, so $f^{(n)}(0)=1$ for all $n\ge0$.
Substituting in \eqref{eq:mm-maclaurin} gives
\[
 e^x=\sum_{n=0}^{\infty}\frac{x^n}{n!}
 =1+x+\frac{x^2}{2!}+\frac{x^3}{3!}+\cdots.
\]
This series converges to $e^x$ for every real $x$, so the approximation can
be replaced by equality. The series can also be used as a definition of the
exponential function.
\end{example}

\subsection{Similar examples: sine and cosine}
The corresponding expansions are
\begin{align*}
 \sin x
 &=\sum_{n=0}^{\infty}\frac{(-1)^n x^{2n+1}}{(2n+1)!}
 =x-\frac{x^3}{3!}+\frac{x^5}{5!}-\cdots,\\[4pt]
 \cos x
 &=\sum_{n=0}^{\infty}\frac{(-1)^n x^{2n}}{(2n)!}
 =1-\frac{x^2}{2!}+\frac{x^4}{4!}-\cdots.
\end{align*}
Both equalities hold for all real $x$. In the lecture's expanded forms,
$3!=6$, $5!=120$ and $4!=24$.

\begin{example}
\begin{align*}
f(x)&=log(1+x), |x|<1.\\
f(0)&=log1=0\\
f'(x)&=\frac{1}{1+x} \rightarrow f'(0)=1\\
f''(x)&=-\frac{1}{(1+x)^{2}}=-(1+x)^{-2} \rightarrow f''(0)=-1\\
f'''(x)&=2(1+   x)^{-3}\rightarrow f'''(0)=2\\
....\\
f^{(n)}(x)&=(-1)^{n-1}(n-1)!\\
\end{align*}
Maclaurin's Series:
\begin{align*}
log(1+x)&\approx\sum_{n=0}^{\infty} (-1)^{n-1} \frac{(n-1)!}{n!}x^{n}\\
\Rightarrow \sum_{n=1}^{\infty} (-1)^{n-1}\frac{x^{n}}{n}&=x-\frac{x^2}{2}+\frac{x^3}{3}-...
\end{align*}
\end{example}

\begin{example}[Binomial Expansion]
\[
f(x)=(1+x)^\alpha,\qquad \alpha\in\mathbb{R}.
\]
For $|x|<1$, its Maclaurin series is
\begin{align*}
(1+x)^\alpha
&=1+\sum_{n=1}^{\infty}
  \frac{\alpha(\alpha-1)\cdots(\alpha-n+1)}{n!}x^n \\
&=\sum_{n=0}^{\infty}\binom{\alpha}{n}x^n.
\end{align*}
Here $\binom{\alpha}{0}=1$.

\textbf{NB:} If $\alpha$ is a non-negative integer, the series terminates
and the identity holds for all real $x$. Otherwise, convergence at
$x=1$ and $x=-1$ must be checked separately.
\end{example}

\begin{example}
Maclaurin Series o
\end{example}

